\section{Related Work}

\label{sec:related}
% ══════════════════════════════════════════════════════════════════════════════

\subsection{Species Identification with Deep Learning}

The past decade has seen dramatic advances in automated species
identification. iNaturalist's computer vision system classifies over
76{,}000 species \citep{van2021benchmarking}. eBird's Merlin app
achieves high accuracy for North American birds using audio and image
inputs \citep{kahl2021birdnet}. PlantNet identifies over 38{,}000 plant
species from photographs \citep{goeau2022overview}. These systems rely
on extensive labeled datasets, with performance degrading sharply for
species with fewer than 100 training images.

\subsection{Few-Shot Learning}

Few-shot learning aims to classify novel classes from $K$ labeled
examples per class (a ``$K$-shot'' setting). Metric-based approaches
learn an embedding space where classification reduces to nearest-neighbor
comparison: matching networks \citep{vinyals2016matching}, prototypical
networks \citep{snell2017prototypical}, and relation networks
\citep{sung2018learning}. Optimization-based methods learn
initialization points for rapid adaptation: MAML \citep{finn2017model}
and its variants. More recently, large pre-trained models with
in-context learning have shown strong few-shot capabilities
\citep{brown2020language}.

\subsection{Hierarchical Classification}

Taxonomic hierarchy has been exploited in traditional species
classification through hierarchical softmax \citep{deng2014large},
label embedding with taxonomy trees \citep{akata2015label}, and
curriculum learning from coarse to fine \citep{bengio2009curriculum}.
\citet{chang2021your} incorporate hierarchical labels into contrastive
learning. However, hierarchical few-shot learning remains
underexplored, with only \citet{li2019large} and
\citet{liu2020many} addressing it in non-ecological contexts.

\subsection{Fine-Grained Recognition}

Fine-grained visual classification distinguishes subordinate categories
(species, cultivars, model variants). Key approaches include bilinear
pooling \citep{lin2015bilinear}, part-based attention
\citep{zheng2017learning}, and feature recalibration
\citep{hu2018squeeze}. Vision transformers (ViTs) have become dominant
due to their ability to capture global context
\citep{dosovitskiy2021image, he2022transfg}. Few-shot fine-grained
recognition remains challenging due to the high intra-class and low
inter-class variance \citep{wei2019piecewise}.


% ══════════════════════════════════════════════════════════════════════════════
